\section*{Chapter \ref{ch:b3:transfers-and-incentives} --- Transfers, Incentives and Redistribution}

\begin{solution}{exo:b3:transfers-and-incentives:1}
It peaks at 40.7~\%, from about 38{,}500 to 55{,}768 dollars, where the earned income credit
is withdrawn at 21.06~\%, \omterm{def:b3:efficiency-cost-of-taxation:payroll-tax}{payroll tax} takes 7.65~\% and income tax 12~\%. Above 55{,}768 the
credit is exhausted and the rate falls to 19.65~\%: \omterm{def:b3:efficiency-cost-of-taxation:payroll-tax}{payroll tax} and the 12~\% bracket.
% figdata: mtr, mtr_stats
\end{solution}

\begin{solution}{exo:b3:transfers-and-incentives:2}
From 10{,}000 to 30{,}000: $(36{,}322 - 14{,}360)/20{,}000 = 1.10$ dollars per dollar. From 30{,}000 to
50{,}000: $(48{,}349 - 36{,}322)/20{,}000 = 0.60$. % figdata: mtr_points
\end{solution}

\begin{solution}{exo:b3:transfers-and-incentives:3}
Break-even $8{,}000/0.4 = 20{,}000$. Transfers: $8{,}000 - 2{,}000 = 6{,}000$ at 5{,}000 and $8{,}000 - 6{,}000 = 2{,}000$ at
15{,}000.
\end{solution}

\begin{solution}{exo:b3:transfers-and-incentives:4}
Recipients, randomly assigned, working as many hours as the controls, or more;
no fall in participation among those with low earnings; or hours that fall as much
among households above the break-even income, whom the programme does not reach.
\end{solution}

\begin{solution}{exo:b3:transfers-and-incentives:5}
At 8{,}000, net income is $8{,}000 + 9{,}000 + 1{,}200 = 18{,}200$: the \omterm{def:b3:transfers-and-incentives:ptr}{participation tax rate} is
$1 - (18{,}200 - 10{,}200)/8{,}000 = 0$. At 15{,}000, net income is $15{,}000 - 600 + 1{,}200 = 15{,}600$: the rate is
$1 - 5{,}400/15{,}000 = 64~\%$. % figdata: brenland_stats
\end{solution}

\begin{solution}{exo:b3:transfers-and-incentives:6}
$0.8/(0.8 + 2 \times 0.25) = 61.5~\%$; with $e = 0.5$, $0.8/(0.8 + 1) = 44.4~\%$.
\end{solution}

\begin{solution}{exo:b3:transfers-and-incentives:7}
The phase-out range becomes longer. Those already in it face a lower marginal rate
and may work more hours; those just above the old end of the range now receive
the credit and face a higher rate, and may work fewer. Some people now find work
worth taking. The cost rises: more households receive the credit, and those already
receiving it receive more.
\end{solution}

\begin{solution}{exo:b3:transfers-and-incentives:8}
The tested solution is \code{ex08_one_child.py}, in the chapter's \code{exercises/}
folder. The one-child credit phases in at 34.00~\% and out at 15.98~\%. The marginal rate
is $-41.35$~\% over most of the first 10{,}000 dollars, 7.65~\% from 15{,}000 to 20{,}000, and highest,
35.63~\%, around 38{,}700 dollars, against 40.7~\% for two children; above the end of the
phase-out it is 19.65~\% for both. % output: ex08_one_child
\end{solution}

\begin{solution}{exo:b3:transfers-and-incentives:9}
No change in participation when credits for mothers expand, relative to women
without children whom the change does not reach; no bunching of earnings where the
credit is largest; or participation that moves as much with the business cycle in
both groups but not with the credit.
\end{solution}

\begin{solution}{exo:b3:transfers-and-incentives:10}
The proof is in the text. $\partial\tau/\partial g = -ae/(1 - g + ae)^2 < 0$; $\tau$ falls in $a$ and in $e$ because both
appear only in the denominator. The taxpayer chose $z$ to maximise utility, so a
small change in $z$ has no first-order effect on her utility (the envelope theorem):
only the mechanical loss counts, and the behavioural change matters only through
revenue.
\end{solution}

\begin{solution}{exo:b3:transfers-and-incentives:11}
(a) The \omterm{def:b3:transfers-and-incentives:ptr}{participation tax rate} is $G/10{,}000$, so $100G/10{,}000 = G/100$ people do not work, and the
budget is $G \times G/100 = 200{,}000$: $G = 4{,}472$, a rate of 44.7~\%, so 55 work and 45 receive 4{,}472.
(b) Any credit makes the \omterm{def:b3:transfers-and-incentives:ptr}{participation tax rate} negative, so all 100 work and each
receives $200{,}000/100 = 2{,}000$. The credit buys participation; the guarantee pays more to
fewer people, and makes them the ones who do not work.
\end{solution}

\begin{solution}{exo:b3:transfers-and-incentives:12}
The tested solution is \code{ex12_gini.py}, in the chapter's \code{exercises/} folder.
Gini 0.290 for gross income, 0.257 after tax, 0.223 for disposable income: of the
reduction of 0.067, taxes account for 0.033 (49~\%) and transfers for 0.034.
% output: ex12_gini
\end{solution}

\begin{solution}{pb:b3:transfers-and-incentives:brenland}
% figdata: brenland_stats, brenland_budget
\textbf{1.} Support $6{,}000 + 2 \times 1{,}500 = 9{,}000$; with the child benefit of 1{,}200, net income 10{,}200.

\textbf{2.} At 11{,}990: $11{,}990 + 9{,}000 + 1{,}200 = 22{,}190$ (no tax below 12{,}000). At 12{,}000: $12{,}000 + 1{,}200 = 13{,}200$.

\textbf{3.} She loses 8{,}990 for 10 dalens more: a marginal rate of about 90{,}000~\%.

\textbf{4.} Zero: no tax, and the support does not change.

\textbf{5.} 151 households.

\textbf{6.} $6{,}000 + 1{,}500(n - 1)$ below 12{,}000, plus $1{,}200\max(n - 2, 0)$: the rules reproduce the column
for every household.

\textbf{7.} Above 12{,}000 her net income is $g - 0.2(g - 12{,}000) + 1{,}200$; it equals 22{,}190 at $g = 23{,}237.5$,
23{,}240 on a 10-dalen grid.

\textbf{8.} $23{,}240/11 = 2{,}113$ hours, more than a full-time year.

\textbf{9.} 64~\% (exercise 5).

\textbf{10.} 11{,}990 gives a net income of 22{,}190; 15{,}000 gives 15{,}600, for 3{,}010 more
dalens of work.

\textbf{11.} A pile-up of households just below 12{,}000 and a gap just above it.

\textbf{12.} The support rises with $n$, so the amount lost at the line, and the income
needed to make it up, rise with household size.

\textbf{13.} $9{,}000/0.5 = 18{,}000$.

\textbf{14.} Support $9{,}000 - 7{,}500 = 1{,}500$; credit $\min(0.2 \times 15{,}000, 2{,}000) = 2{,}000$; net $15{,}000 - 600 + 1{,}500 + 2{,}000 + 1{,}200 =
19{,}100$.

\textbf{15.} $1 - (19{,}100 - 10{,}200)/15{,}000 = 40.7~\%$.

\textbf{16.} Before, 0; after, net income at 8{,}000 is $8{,}000 + 5{,}000 + 1{,}600 + 1{,}200 = 15{,}800$, a rate of
$1 - 5{,}600/8{,}000 = 30~\%$. The reform raises the rate on small earnings to lower it further up.

\textbf{17.} Working households below 12{,}000 whose support is now partly withdrawn:
a single person earning 10{,}000 received 6{,}000; after the reform $1{,}000 + 2{,}000 = 3{,}000$.

\textbf{18.} 859.7 dalens per household, 1{,}719{,}355 dalens for the 2{,}000, before any response.

\textbf{19.} 1{,}210 gain and 73 lose; the rest are unaffected.

\textbf{20.} 0.290.

\textbf{21.} 0.223 before and 0.211 after.

\textbf{22.} By 0.012, 5.5~\% of the disposable-income Gini.

\textbf{23.} It holds gross incomes fixed. If people work more after the reform, gross
incomes at the bottom rise, and both the Gini and the cost change.

\textbf{24.} More work by parents now facing 41~\% instead of 64~\% lowers the cost (less support, more
tax); less work by those whose small earnings now face 30~\% instead of 0 raises it;
households may also bunch where the credit is largest.

\textbf{25.} The \omterm{def:b3:transfers-and-incentives:ptr}{participation tax rate} of a parent of two at 15{,}000 dalens falls from 64~\%
to 40.7~\%, and the reform costs 860 dalens per household before any response.
\end{solution}
